% Relatório parcial: questões (a), (b) e (c). Compilador: pdfLaTeX.
% Ajustes de layout: fonte do corpo abaixo; tabelas em small; margens em geometry.
\documentclass[11pt,a4paper]{article}
\usepackage[utf8]{inputenc}
\usepackage[T1]{fontenc}
\usepackage[brazil]{babel}
\usepackage{lmodern,textcomp}
\usepackage[margin=2cm]{geometry}
\usepackage{amsmath,amssymb}
\usepackage{longtable,booktabs,array,calc}
\usepackage{microtype,etoolbox}
\setcounter{secnumdepth}{-1}
\setlength{\parindent}{0pt}
\setlength{\parskip}{4pt}
\linespread{1.03}
\setlength{\tabcolsep}{4pt}
\renewcommand{\arraystretch}{1.12}
\setlength{\emergencystretch}{2em}
\AtBeginEnvironment{longtable}{\small}
\providecommand{\tightlist}{\setlength{\itemsep}{0pt}\setlength{\parskip}{0pt}}
\date{}

\begin{document}

\section{Relatório de Diagnóstico --- Mercatto
Varejo}\label{relatuxf3rio-de-diagnuxf3stico-mercatto-varejo}

\textbf{Mineração de Dados · IFMT · Entrega 1}\\
\textbf{Equipe:} Matheus Henrique Moreira Louro e Lucas Ribeiro
Luchesi\\
\textbf{Ambiente:} Docker, escala 1,0, semente 20262

\subsection{1. Objetivo e método}\label{objetivo-e-muxe9todo}

Diagnosticar as fontes MySQL legado, PostgreSQL do e-commerce, CSV do
SAC e NDJSON de avaliações antes da integração. As medições abrangem as
bases completas, sem alteração das origens, e complementam o perfilador
fornecido. O objetivo do projeto é prever a recompra entre 01/07/2025 e
28/09/2025 usando o comportamento até 30/06/2025.

\subsection{2. Questão (a) --- Chaves de junção e
cobertura}\label{questuxe3o-a-chaves-de-junuxe7uxe3o-e-cobertura}

\textbf{Cobertura = 100 \ensuremath{\times} registros com correspondência / registros
examinados.} Cada registro é contado uma vez, mesmo quando sua chave
encontra várias ocorrências. Ausências e duplicatas integram o
denominador; valores vazios não são pareados. CPF e telefone foram
comparados somente por dígitos; e-mail, em minúsculas e sem espaços
externos; nomes, sem acentos, em maiúsculas e com espaços uniformizados.

\subsubsection{2.1. Clientes entre os
sistemas}\label{clientes-entre-os-sistemas}

\begin{longtable}[]{@{}
  >{\raggedright\arraybackslash}p{(\linewidth - 6\tabcolsep) * \real{0.2500}}
  >{\raggedleft\arraybackslash}p{(\linewidth - 6\tabcolsep) * \real{0.1800}}
  >{\raggedleft\arraybackslash}p{(\linewidth - 6\tabcolsep) * \real{0.2850}}
  >{\raggedleft\arraybackslash}p{(\linewidth - 6\tabcolsep) * \real{0.2850}}@{}}
\toprule\noalign{}
\begin{minipage}[b]{\linewidth}\raggedright
Candidato
\end{minipage} & \begin{minipage}[b]{\linewidth}\raggedleft
Valores distintos compartilhados
\end{minipage} & \begin{minipage}[b]{\linewidth}\raggedleft
Legado: cobertos / 47.250
\end{minipage} & \begin{minipage}[b]{\linewidth}\raggedleft
E-commerce: cobertos / 85.000
\end{minipage} \\
\midrule\noalign{}
\endhead
\bottomrule\noalign{}
\endlastfoot
\texttt{CPF} \ensuremath{\leftrightarrow} \texttt{tax\_id} & 12.605 & 13.048 (27,61\%) & 12.606
(14,83\%) \\
\texttt{NOME} \ensuremath{\leftrightarrow} \texttt{full\_name} & 22.869 & 28.204 (59,69\%) & 31.360
(36,89\%) \\
\texttt{EMAIL} \ensuremath{\leftrightarrow} \texttt{email} & 0 & 0 (0,00\%) & 0 (0,00\%) \\
\texttt{TELEFONE} \ensuremath{\leftrightarrow} \texttt{phone} & 0 & 0 (0,00\%) & 0 (0,00\%) \\
\end{longtable}

O \textbf{CPF é o principal candidato}, mas não garante identidade nem
unicidade: 1.427 valores se repetem no legado e 3 no e-commerce; somente
12.161 dos 12.605 CPFs compartilhados são únicos em ambos. Entre os
compartilhados, 1.135 (9,00\%) têm dígito verificador inválido;
excluí-los eliminaria correspondências observáveis. O nome apresenta
maior cobertura, porém há 6.758 nomes repetidos no legado e 15.477 no
e-commerce, exigindo critérios adicionais para evitar unir pessoas
distintas. E-mail e telefone não ligam os cadastros nesta base. Os IDs
internos dos dois sistemas também não comprovam identidade entre fontes.

\subsubsection{2.2. Arquivos para pedidos, clientes e
produtos}\label{arquivos-para-pedidos-clientes-e-produtos}

\begin{longtable}[]{@{}
  >{\raggedright\arraybackslash}p{(\linewidth - 4\tabcolsep) * \real{0.5800}}
  >{\raggedleft\arraybackslash}p{(\linewidth - 4\tabcolsep) * \real{0.2800}}
  >{\raggedleft\arraybackslash}p{(\linewidth - 4\tabcolsep) * \real{0.1400}}@{}}
\toprule\noalign{}
\begin{minipage}[b]{\linewidth}\raggedright
Candidato e destino
\end{minipage} & \begin{minipage}[b]{\linewidth}\raggedleft
Cobertos / examinados
\end{minipage} & \begin{minipage}[b]{\linewidth}\raggedleft
Cobertura
\end{minipage} \\
\midrule\noalign{}
\endhead
\bottomrule\noalign{}
\endlastfoot
SAC: \texttt{(origem\_sistema,\ id\_pedido\_origem)} \ensuremath{\rightarrow} pedido na origem
indicada & 96.140 / 96.140 & 100,00\% \\
SAC: \texttt{cpf\_cliente} \ensuremath{\rightarrow} qualquer cadastro & 90.155 / 96.140 &
93,77\% \\
SAC: \texttt{cpf\_cliente} \ensuremath{\rightarrow} cadastro da origem indicada & 88.535 /
96.140 & 92,09\% \\
JSON: \texttt{order.ref} \ensuremath{\rightarrow} pedido da origem indicada & 160.000 / 160.000
& 100,00\% \\
JSON: \texttt{product.sku} \ensuremath{\rightarrow} produto da origem indicada & 160.000 /
160.000 & 100,00\% \\
\end{longtable}

A chave de pedido deve incluir a origem, pois IDs podem coincidir entre
os bancos. No SAC, \texttt{LEGADO} e \texttt{ECOM} direcionam a busca a
\texttt{ID\_PEDIDO} e \texttt{order\_id}. No JSON,
\texttt{LEG-\textless{}id\textgreater{}} referencia IDs do legado,
\texttt{ECOM-\textless{}id\textgreater{}} referencia pedidos do
e-commerce e \texttt{SKU-\textless{}n\textgreater{}} corresponde a
\texttt{shop.product.sku}. São 40.016 avaliações do legado e 119.984 do
e-commerce, com cobertura de 100\% para pedidos e produtos em cada
grupo.

A ligação do SAC por pedido cobre todas as linhas, inclusive as 5.985
(6,23\%) sem CPF em nenhum cadastro. Nas avaliações, a existência de
pedido e produto foi verificada separadamente; não comprova que o
produto pertence ao pedido. \texttt{protocolo} e \texttt{review\_id}
identificam registros locais, sem contraparte direta nos bancos.

\subsubsection{2.3. Relações internas dos
bancos}\label{relauxe7uxf5es-internas-dos-bancos}

\begin{longtable}[]{@{}
  >{\raggedright\arraybackslash}p{(\linewidth - 4\tabcolsep) * \real{0.4000}}
  >{\raggedleft\arraybackslash}p{(\linewidth - 4\tabcolsep) * \real{0.3000}}
  >{\raggedleft\arraybackslash}p{(\linewidth - 4\tabcolsep) * \real{0.3000}}@{}}
\toprule\noalign{}
\begin{minipage}[b]{\linewidth}\raggedright
Relação e campo da chave
\end{minipage} & \begin{minipage}[b]{\linewidth}\raggedleft
Legado: cobertos / examinados
\end{minipage} & \begin{minipage}[b]{\linewidth}\raggedleft
E-commerce: cobertos / examinados
\end{minipage} \\
\midrule\noalign{}
\endhead
\bottomrule\noalign{}
\endlastfoot
Pedido \ensuremath{\rightarrow} cliente: \texttt{ID\_CLIENTE} / \texttt{customer\_id} & 130.000
/ 130.000 (100\%) & 280.000 / 280.000 (100\%) \\
Item \ensuremath{\rightarrow} pedido: \texttt{ID\_PEDIDO} / \texttt{order\_id} & 391.218 /
397.086 (98,52\%) & 980.658 / 995.367 (98,52\%) \\
Item \ensuremath{\rightarrow} produto: \texttt{ID\_PRODUTO} / \texttt{product\_id} & 397.086 /
397.086 (100\%) & 995.367 / 995.367 (100\%) \\
\end{longtable}

Há 5.868 itens sem pedido no legado e 14.709 no e-commerce (1,48\% em
cada origem). IDs de produto de sistemas diferentes não estabelecem
equivalência comercial.

\subsection{3. Questão (b) --- Heterogeneidade
semântica}\label{questuxe3o-b-heterogeneidade-semuxe2ntica}

Foram identificados três casos em que campos aparentemente equivalentes
representam conceitos diferentes. Variações de idioma, grafia ou
formato, isoladamente, não constituem divergência semântica.

\subsubsection{3.1. Cancelamento e recusa de
pagamento}\label{cancelamento-e-recusa-de-pagamento}

Segundo o esqueleto de conformação fornecido, \texttt{CANCELADO} no
legado agrega cancelamento pelo cliente e recusa de pagamento; o
e-commerce distingue \texttt{canceled} de \texttt{payment\_declined}. As
distribuições confirmam a diferença de abrangência:

\begin{longtable}[]{@{}lrr@{}}
\toprule\noalign{}
Sistema e status & Pedidos / total da origem & Percentual \\
\midrule\noalign{}
\endhead
\bottomrule\noalign{}
\endlastfoot
Legado: \texttt{CANCELADO} & 15.774 / 130.000 & 12,13\% \\
E-commerce: \texttt{canceled} & 22.413 / 280.000 & 8,00\% \\
E-commerce: \texttt{payment\_declined} & 13.990 / 280.000 & 5,00\% \\
\end{longtable}

A soma dos dois estados no e-commerce é 36.403 pedidos (13,00\%).
Comparar apenas \texttt{canceled} com \texttt{CANCELADO} excluiria as
recusas da segunda fonte. Agrupar os estados permite comparação mais
ampla, mas perde a distinção disponível no e-commerce. O motivo
específico não pode ser recuperado no legado.

\subsubsection{3.2. Composição do valor do
pedido}\label{composiuxe7uxe3o-do-valor-do-pedido}

\texttt{VALOR\_TOTAL} do legado inclui frete; \texttt{total\_amount} do
e-commerce corresponde ao valor líquido dos itens. Essa diferença foi
confirmada nos 130.000 pedidos do legado e nos 280.000 do e-commerce
(100\% de cada origem), com resíduo máximo de R\$ 0,00 nas respectivas
fórmulas. Comparar ou somar esses campos diretamente misturaria medidas
distintas. A prova e os critérios de reconciliação estão na questão (c).

\subsubsection{3.3. Abrangência das categorias de
produto}\label{abranguxeancia-das-categorias-de-produto}

O legado tem 27 rótulos distintos, considerando caixa e acentos, para
3.000 produtos; o e-commerce tem 21 nós de categoria (8 raízes e 13
filhos), para 12.000 produtos. ``Categoria'' pode representar um ramo,
uma categoria específica ou uma combinação de ramos:

\begin{itemize}
\tightlist
\item
  \texttt{informatica\ e\ games}: 112 produtos (3,73\% do legado),
  reunindo dois ramos separados no e-commerce.
\item
  \texttt{TV\ E\ SOM}: 111 (3,70\%), reunindo \texttt{Televisores} e
  \texttt{Áudio}, que classificam 793 e 805 produtos no e-commerce.
\item
  \texttt{DIVERSOS}, \texttt{OUTROS} e vazio: 340 (11,33\%), sem ramo
  comercial identificável.
\end{itemize}

Padronizar grafias não resolve essas diferenças de significado. Atribuir
uma categoria combinada a um único ramo descartaria parte da
classificação; separá-la exige informação adicional por produto. Os 27
rótulos, portanto, não equivalem a 27 categorias comerciais.

\subsection{4. Questão (c) --- Reconciliação
financeira}\label{questuxe3o-c-reconciliauxe7uxe3o-financeira}

\textbf{Os campos não medem a mesma coisa:} \texttt{VALOR\_TOTAL} inclui
frete, enquanto \texttt{total\_amount} o exclui. Para verificar, os
itens foram agregados por ID de pedido antes da junção ao cabeçalho,
evitando repetir totais e fretes. Foram examinados todos os pedidos,
independentemente do status; itens órfãos não entram na soma de nenhum
pedido existente.

Para cada pedido \(p\), definiu-se o valor líquido dos itens:

\[L_p = \sum_{i \in I_p} (q_i u_i - d_i^*)\]

Em que \(I_p\) são os itens do pedido, \(q_i\) é a quantidade, \(u_i\) o
preço unitário praticado e \(d_i^*\) o desconto da linha. Foram usados
\texttt{QTD}, \texttt{VL\_UNITARIO} e \texttt{VL\_DESCONTO} no legado;
\texttt{quantity}, \texttt{unit\_price} e \texttt{discount} no
e-commerce.

No legado, o sentinela \texttt{VL\_DESCONTO\ =\ -1} contribuiu como zero
no teste. Essa hipótese é compatível com os totais registrados, mas não
recupera o desconto real ausente. Há 4.056 itens com esse sentinela em
4.003 pedidos (3,08\% dos 130.000); subtrair literalmente \texttt{-1}
acrescentaria R\$ 1,00 por item e faria esses pedidos falharem na
reconciliação. Os fretes foram usados com o sinal original, inclusive
quando negativos.

\begin{longtable}[]{@{}
  >{\raggedright\arraybackslash}p{(\linewidth - 4\tabcolsep) * \real{0.3800}}
  >{\raggedleft\arraybackslash}p{(\linewidth - 4\tabcolsep) * \real{0.3100}}
  >{\raggedleft\arraybackslash}p{(\linewidth - 4\tabcolsep) * \real{0.3100}}@{}}
\toprule\noalign{}
\begin{minipage}[b]{\linewidth}\raggedright
Hipótese para o total registrado
\end{minipage} & \begin{minipage}[b]{\linewidth}\raggedleft
Legado: 130.000 pedidos
\end{minipage} & \begin{minipage}[b]{\linewidth}\raggedleft
E-commerce: 280.000 pedidos
\end{minipage} \\
\midrule\noalign{}
\endhead
\bottomrule\noalign{}
\endlastfoot
\(L_p + \text{frete}_p\) & 130.000 (100,00\%) & 94 (0,03\%) \\
\(L_p\), sem frete & 33 (0,03\%) & 280.000 (100,00\%) \\
\end{longtable}

A tolerância foi R\$ 0,02 por pedido. Nas fórmulas confirmadas, o maior
resíduo absoluto foi \textbf{R\$ 0,00}. As 33 e 94 coincidências com a
hipótese alternativa decorrem de fretes cuja magnitude está dentro da
tolerância. A reconciliação agregada também confirma o resultado:

\begin{center}\small\begin{tabular}{@{}lrr@{}}
\toprule\noalign{}
Medida agregada (R\$) & Legado & E-commerce \\
\midrule\noalign{}
Soma dos totais registrados & 281.084.957,13 & 734.859.099,95 \\
Soma dos itens líquidos & 275.324.024,56 & 734.859.099,95 \\
Soma dos fretes & 5.760.932,57 & 10.991.760,38 \\
\bottomrule\end{tabular}\end{center}

No legado, \textbf{281.084.957,13 = 275.324.024,56 + 5.760.932,57}. No
e-commerce, a soma dos totais já é igual à soma dos itens líquidos;
incluir frete exigiria acrescentar R\$ 10.991.760,38. Portanto, uma
medida comparável de total com frete deve manter \texttt{VALOR\_TOTAL}
no legado e adicionar \texttt{freight\_amount} a \texttt{total\_amount}
no e-commerce, preservando a origem das medidas.

\subsection{5. Questão (d) --- Anomalias e limites do
diagnóstico}\label{questuxe3o-d-anomalias-e-limites-do-diagnuxf3stico}

O diagnóstico das fontes completas permitiu quantificar anomalias locais estruturais e semânticas que afetam a qualidade dos dados antes da integração:

\begin{itemize}
\tightlist
\item
  \textbf{Valores sentinela:} No legado, 1.025 dos 130.000 pedidos (0,79\%) registram a data de entrega como \texttt{1900-01-01}. Esse artifício para contornar a ausência de data realiza um deslocamento nas médias cronológicas.
\item
  \textbf{Inconsistências cronológicas e financeiras:} Entregas registradas antes do próprio pedido ocorrem em 775 pedidos do legado (0,60\%) e 1.683 do e-commerce (0,60\%). Além disso, há valores de frete negativos em 362 pedidos do legado (0,28\%) e 812 do e-commerce (0,29\%). Estes representam defeitos sistêmicos ou cancelamentos compensatórios que exigem investigação.
\item
  \textbf{Referências órfãs:} A integridade referencial está quebrada em ambas as origens. Há 5.868 itens sem pedido correspondente no legado (1,48\% de 397.086 itens) e 14.709 no e-commerce (1,48\% de 995.367).
\item
  \textbf{Qualidade e tipagem em arquivos:} O SAC (96.140 chamados) tem 2.830 notas de satisfação fora do domínio esperado (1 a 5) e 1.377 registros contêm quebras de linha internas, exigindo tratamento de conversão e formatação. No JSON (160.000 avaliações), 4.831 registros (3,02\%) omitem o campo \texttt{rating}, e dentre os presentes, 16.108 (10,38\%) são formatados como texto (\texttt{string}) em vez de inteiro.
\end{itemize}

\textbf{Limites do diagnóstico e problemas revelados apenas na integração:}\\
As anomalias supracitadas foram detectadas por consultas e perfilamentos locais independentes. No entanto, os defeitos semânticos mais graves dependem do resultado do cruzamento das bases e da aplicação de regras de negócio integradas.
A divergência principal é a composição do valor dos pedidos: conforme provado na questão (c), o legado embute o frete no total, mas o e-commerce o isola. Um diagnóstico unifonte não detecta problemas nesses valores, pois eles são válidos localmente. Integrá-los sem uma regra de equiparação causaria viés analítico ao superestimar a receita de clientes antigos. Da mesma forma, o tratamento de identidades duplicadas e CPFs incorretos (presentes igualmente nas duas bases) dita a sobrevivência do registro. Se as regras de integração aglutinarem equivocadamente pessoas distintas, o modelo preditivo será treinado com históricos contaminados e hiperinflacionados.

\subsection{6. Questão (e) --- Grão das
tabelas-fato}\label{questuxe3o-e-gruxe3o-das-tabelas-fato}

O modelo dimensional estabelece as seguintes unidades de análise para preservar a rastreabilidade e evitar o produto cartesiano ao agregar as métricas:

\begin{itemize}
\tightlist
\item
  Uma linha de \texttt{dw.fato\_pedido} representa um único pedido comercializado em qualquer um dos sistemas de origem.
\item
  Uma linha de \texttt{dw.fato\_item\_pedido} representa um produto unitário faturado dentro de um pedido específico.
\item
  Uma linha de \texttt{dw.fato\_atendimento} representa um registro de atendimento (chamado) aberto no SAC.
\item
  Uma linha de \texttt{dw.fato\_avaliacao} representa uma avaliação de produto enviada por um cliente ao marketplace.
\end{itemize}

A chave primária local dos pedidos (\texttt{ID\_PEDIDO} no legado e \texttt{order\_id} no e-commerce) é insuficiente para garantir a unicidade no ambiente integrado. Os sistemas geradores operam de forma isolada e emitem sequências numéricas idênticas. Devido a essa colisão inevitável de identificadores, a origem compõe a identificação obrigatória no Data Warehouse (através da chave composta por \texttt{id\_origem} e \texttt{sistema\_origem}), assegurando que um pedido oriundo do legado e um do e-commerce com o mesmo número sejam corretamente apartados e reconhecidos como eventos distintos.

\subsection{Referências e
reprodução}\label{referuxeancias-e-reproduuxe7uxe3o}

Projeto Prático ETL Mercatto, seção 3.1; Dicionário de Dados e esqueleto
de transformação do laboratório. Medições:
\texttt{diagnostico\_chaves.py} e \texttt{diagnostico\_semantica.py}, na
pasta \texttt{etl}; resultados em arquivos JSON homônimos, na pasta
\texttt{saida}.

\end{document}
